\documentclass[12pt, openany]{book}
\usepackage[utf8]{inputenc}
\usepackage[spanish]{babel}
\usepackage{amsmath}
\usepackage{amssymb}
\usepackage{geometry}
\usepackage{listings}
\usepackage{xcolor}
\usepackage{hyperref}
\usepackage{setspace}

\geometry{a4paper, margin=1.2in}
\onehalfspacing

\definecolor{codegreen}{rgb}{0,0.6,0}
\definecolor{codegray}{rgb}{0.5,0.5,0.5}
\definecolor{codepurple}{rgb}{0.58,0,0.82}
\definecolor{backcolour}{rgb}{0.95,0.95,0.92}

\lstdefinestyle{mystyle}{
    backgroundcolor=\color{backcolour},   
    commentstyle=\color{codegreen},
    keywordstyle=\color{magenta},
    numberstyle=\tiny\color{codegray},
    stringstyle=\color{codepurple},
    basicstyle=\ttfamily\small,
    breakatwhitespace=false,         
    breaklines=true,                 
    captionpos=b,                    
    keepspaces=true,                 
    numbers=left,                    
    numbersep=5pt,                  
    showspaces=false,                
    showstringspaces=false,
    showtabs=false,                  
    tabsize=2
}
\lstset{style=mystyle}

\title{\Huge \textbf{Cosmic-Generator} \\ \vspace{0.5cm} \Huge \textbf{Manual de Usuario - Volumen 2} \\ \vspace{1cm} \Large \textit{La Enciclopedia Cósmica: Teoría, Código y Caos}}
\author{Cooperativa Central: Rebeldía Cósmica}
\date{\today}

\begin{document}
\maketitle
\tableofcontents
\newpage
\part{Fundamentos Teóricos: Sistemas Complejos y Redes Neuronales}


\chapter{Introducción a los Sistemas Complejos}
La teoría de los sistemas complejos estudia cómo las interacciones entre partes simples dan lugar a comportamientos colectivos emergentes, impredecibles a partir de las propiedades individuales. En el contexto de \textbf{Cosmic-Generator}, empleamos ecuaciones diferenciales parciales (EDPs) no lineales para modelar dinámicas de fluidos, fuego, ecosistemas bacterianos y fenómenos cuánticos.

\section{El Surgimiento del Caos Determinista}
El caos determinista es un fenómeno en el cual sistemas dinámicos que están gobernados por leyes matemáticas estrictas y sin componentes aleatorios exhiben comportamientos impredecibles a largo plazo debido a su extrema sensibilidad a las condiciones iniciales. Este efecto, popularmente conocido como el "Efecto Mariposa", fue documentado rigurosamente por Edward Lorenz en 1963 mientras modelaba la convección atmosférica. 

En nuestro motor, implementamos precisamente el Atractor de Lorenz como el latido invisible que empuja las capas ocultas de nuestra red neuronal, logrando que el arte generativo nunca se repita exactamente igual, a pesar de estar estrictamente acoplado al espectro de audio. Las ecuaciones de Lorenz son:

\begin{align}
\frac{dx}{dt} &= \sigma(y - x) \\
\frac{dy}{dt} &= x(\rho - z) - y \\
\frac{dz}{dt} &= xy - \beta z
\end{align}

\section{Reacción-Difusión y la Morfogénesis de Turing}
Alan Turing postuló en 1952 que un sistema de sustancias químicas reaccionando y difundiendo en un espacio podía generar patrones espacialmente periódicos a partir de un estado homogéneo debido a la inestabilidad. En Cosmic-Generator usamos el modelo Gray-Scott para materializar esta visión, inyectando audio para desestabilizar continuamente el medio.



\chapter{Introducción a los Sistemas Complejos}
La teoría de los sistemas complejos estudia cómo las interacciones entre partes simples dan lugar a comportamientos colectivos emergentes, impredecibles a partir de las propiedades individuales. En el contexto de \textbf{Cosmic-Generator}, empleamos ecuaciones diferenciales parciales (EDPs) no lineales para modelar dinámicas de fluidos, fuego, ecosistemas bacterianos y fenómenos cuánticos.

\section{El Surgimiento del Caos Determinista}
El caos determinista es un fenómeno en el cual sistemas dinámicos que están gobernados por leyes matemáticas estrictas y sin componentes aleatorios exhiben comportamientos impredecibles a largo plazo debido a su extrema sensibilidad a las condiciones iniciales. Este efecto, popularmente conocido como el "Efecto Mariposa", fue documentado rigurosamente por Edward Lorenz en 1963 mientras modelaba la convección atmosférica. 

En nuestro motor, implementamos precisamente el Atractor de Lorenz como el latido invisible que empuja las capas ocultas de nuestra red neuronal, logrando que el arte generativo nunca se repita exactamente igual, a pesar de estar estrictamente acoplado al espectro de audio. Las ecuaciones de Lorenz son:

\begin{align}
\frac{dx}{dt} &= \sigma(y - x) \\
\frac{dy}{dt} &= x(\rho - z) - y \\
\frac{dz}{dt} &= xy - \beta z
\end{align}

\section{Reacción-Difusión y la Morfogénesis de Turing}
Alan Turing postuló en 1952 que un sistema de sustancias químicas reaccionando y difundiendo en un espacio podía generar patrones espacialmente periódicos a partir de un estado homogéneo debido a la inestabilidad. En Cosmic-Generator usamos el modelo Gray-Scott para materializar esta visión, inyectando audio para desestabilizar continuamente el medio.



\chapter{Introducción a los Sistemas Complejos}
La teoría de los sistemas complejos estudia cómo las interacciones entre partes simples dan lugar a comportamientos colectivos emergentes, impredecibles a partir de las propiedades individuales. En el contexto de \textbf{Cosmic-Generator}, empleamos ecuaciones diferenciales parciales (EDPs) no lineales para modelar dinámicas de fluidos, fuego, ecosistemas bacterianos y fenómenos cuánticos.

\section{El Surgimiento del Caos Determinista}
El caos determinista es un fenómeno en el cual sistemas dinámicos que están gobernados por leyes matemáticas estrictas y sin componentes aleatorios exhiben comportamientos impredecibles a largo plazo debido a su extrema sensibilidad a las condiciones iniciales. Este efecto, popularmente conocido como el "Efecto Mariposa", fue documentado rigurosamente por Edward Lorenz en 1963 mientras modelaba la convección atmosférica. 

En nuestro motor, implementamos precisamente el Atractor de Lorenz como el latido invisible que empuja las capas ocultas de nuestra red neuronal, logrando que el arte generativo nunca se repita exactamente igual, a pesar de estar estrictamente acoplado al espectro de audio. Las ecuaciones de Lorenz son:

\begin{align}
\frac{dx}{dt} &= \sigma(y - x) \\
\frac{dy}{dt} &= x(\rho - z) - y \\
\frac{dz}{dt} &= xy - \beta z
\end{align}

\section{Reacción-Difusión y la Morfogénesis de Turing}
Alan Turing postuló en 1952 que un sistema de sustancias químicas reaccionando y difundiendo en un espacio podía generar patrones espacialmente periódicos a partir de un estado homogéneo debido a la inestabilidad. En Cosmic-Generator usamos el modelo Gray-Scott para materializar esta visión, inyectando audio para desestabilizar continuamente el medio.



\chapter{Introducción a los Sistemas Complejos}
La teoría de los sistemas complejos estudia cómo las interacciones entre partes simples dan lugar a comportamientos colectivos emergentes, impredecibles a partir de las propiedades individuales. En el contexto de \textbf{Cosmic-Generator}, empleamos ecuaciones diferenciales parciales (EDPs) no lineales para modelar dinámicas de fluidos, fuego, ecosistemas bacterianos y fenómenos cuánticos.

\section{El Surgimiento del Caos Determinista}
El caos determinista es un fenómeno en el cual sistemas dinámicos que están gobernados por leyes matemáticas estrictas y sin componentes aleatorios exhiben comportamientos impredecibles a largo plazo debido a su extrema sensibilidad a las condiciones iniciales. Este efecto, popularmente conocido como el "Efecto Mariposa", fue documentado rigurosamente por Edward Lorenz en 1963 mientras modelaba la convección atmosférica. 

En nuestro motor, implementamos precisamente el Atractor de Lorenz como el latido invisible que empuja las capas ocultas de nuestra red neuronal, logrando que el arte generativo nunca se repita exactamente igual, a pesar de estar estrictamente acoplado al espectro de audio. Las ecuaciones de Lorenz son:

\begin{align}
\frac{dx}{dt} &= \sigma(y - x) \\
\frac{dy}{dt} &= x(\rho - z) - y \\
\frac{dz}{dt} &= xy - \beta z
\end{align}

\section{Reacción-Difusión y la Morfogénesis de Turing}
Alan Turing postuló en 1952 que un sistema de sustancias químicas reaccionando y difundiendo en un espacio podía generar patrones espacialmente periódicos a partir de un estado homogéneo debido a la inestabilidad. En Cosmic-Generator usamos el modelo Gray-Scott para materializar esta visión, inyectando audio para desestabilizar continuamente el medio.



\chapter{Redes Neuronales Artificiales y CPPNs}
Las Redes Neuronales Artificiales (ANNs) nacen de la inspiración biológica del cerebro humano. Sin embargo, en \textbf{Cosmic-Generator}, no utilizamos ANNs tradicionales para clasificación o regresión, sino \textbf{Compositional Pattern Producing Networks (CPPNs)}.

\section{Arquitectura del Modo Neural (CPPN)}
Una CPPN es una red neuronal generativa que mapea coordenadas espaciales $(x, y)$ a un valor de color (o densidad) en ese píxel. A diferencia de las redes convolucionales o perceptrones multicapa clásicos que usan la misma función de activación (ej. ReLU) en todas las neuronas, las CPPNs utilizan un cóctel de funciones matemáticas heterogéneas (seno, coseno, hiperbólicas, gaussianas). Esto permite que la red componga patrones que contienen simetría (debido al coseno), repetición (debido al seno) y estructuras centradas (debido a la gaussiana).

\section{El Modo Neural en Cosmic-Generator}
Nuestro motor inicializa una CPPN con pesos completamente aleatorios. No hay fase de entrenamiento (Backpropagation). En lugar de ello, la red recibe como entrada las coordenadas normalizadas del lienzo, el radio respecto al centro, y el estado actual de un atractor de Lorenz tridimensional. La salida es una matriz de intensidades que se mapea a través de la \textit{Cromestesia} (la paleta de colores reactiva al tono de la música).



\chapter{Redes Neuronales Artificiales y CPPNs}
Las Redes Neuronales Artificiales (ANNs) nacen de la inspiración biológica del cerebro humano. Sin embargo, en \textbf{Cosmic-Generator}, no utilizamos ANNs tradicionales para clasificación o regresión, sino \textbf{Compositional Pattern Producing Networks (CPPNs)}.

\section{Arquitectura del Modo Neural (CPPN)}
Una CPPN es una red neuronal generativa que mapea coordenadas espaciales $(x, y)$ a un valor de color (o densidad) en ese píxel. A diferencia de las redes convolucionales o perceptrones multicapa clásicos que usan la misma función de activación (ej. ReLU) en todas las neuronas, las CPPNs utilizan un cóctel de funciones matemáticas heterogéneas (seno, coseno, hiperbólicas, gaussianas). Esto permite que la red componga patrones que contienen simetría (debido al coseno), repetición (debido al seno) y estructuras centradas (debido a la gaussiana).

\section{El Modo Neural en Cosmic-Generator}
Nuestro motor inicializa una CPPN con pesos completamente aleatorios. No hay fase de entrenamiento (Backpropagation). En lugar de ello, la red recibe como entrada las coordenadas normalizadas del lienzo, el radio respecto al centro, y el estado actual de un atractor de Lorenz tridimensional. La salida es una matriz de intensidades que se mapea a través de la \textit{Cromestesia} (la paleta de colores reactiva al tono de la música).



\chapter{Redes Neuronales Artificiales y CPPNs}
Las Redes Neuronales Artificiales (ANNs) nacen de la inspiración biológica del cerebro humano. Sin embargo, en \textbf{Cosmic-Generator}, no utilizamos ANNs tradicionales para clasificación o regresión, sino \textbf{Compositional Pattern Producing Networks (CPPNs)}.

\section{Arquitectura del Modo Neural (CPPN)}
Una CPPN es una red neuronal generativa que mapea coordenadas espaciales $(x, y)$ a un valor de color (o densidad) en ese píxel. A diferencia de las redes convolucionales o perceptrones multicapa clásicos que usan la misma función de activación (ej. ReLU) en todas las neuronas, las CPPNs utilizan un cóctel de funciones matemáticas heterogéneas (seno, coseno, hiperbólicas, gaussianas). Esto permite que la red componga patrones que contienen simetría (debido al coseno), repetición (debido al seno) y estructuras centradas (debido a la gaussiana).

\section{El Modo Neural en Cosmic-Generator}
Nuestro motor inicializa una CPPN con pesos completamente aleatorios. No hay fase de entrenamiento (Backpropagation). En lugar de ello, la red recibe como entrada las coordenadas normalizadas del lienzo, el radio respecto al centro, y el estado actual de un atractor de Lorenz tridimensional. La salida es una matriz de intensidades que se mapea a través de la \textit{Cromestesia} (la paleta de colores reactiva al tono de la música).



\chapter{Redes Neuronales Artificiales y CPPNs}
Las Redes Neuronales Artificiales (ANNs) nacen de la inspiración biológica del cerebro humano. Sin embargo, en \textbf{Cosmic-Generator}, no utilizamos ANNs tradicionales para clasificación o regresión, sino \textbf{Compositional Pattern Producing Networks (CPPNs)}.

\section{Arquitectura del Modo Neural (CPPN)}
Una CPPN es una red neuronal generativa que mapea coordenadas espaciales $(x, y)$ a un valor de color (o densidad) en ese píxel. A diferencia de las redes convolucionales o perceptrones multicapa clásicos que usan la misma función de activación (ej. ReLU) en todas las neuronas, las CPPNs utilizan un cóctel de funciones matemáticas heterogéneas (seno, coseno, hiperbólicas, gaussianas). Esto permite que la red componga patrones que contienen simetría (debido al coseno), repetición (debido al seno) y estructuras centradas (debido a la gaussiana).

\section{El Modo Neural en Cosmic-Generator}
Nuestro motor inicializa una CPPN con pesos completamente aleatorios. No hay fase de entrenamiento (Backpropagation). En lugar de ello, la red recibe como entrada las coordenadas normalizadas del lienzo, el radio respecto al centro, y el estado actual de un atractor de Lorenz tridimensional. La salida es una matriz de intensidades que se mapea a través de la \textit{Cromestesia} (la paleta de colores reactiva al tono de la música).


\part{Ingeniería de Código: Discretizaciones y Fronteras}


\chapter{Discretizaciones Espaciales}
Para simular el universo continuo en una cuadrícula de píxeles, debemos discretizar los operadores diferenciales continuos (como el Laplaciano $\nabla^2$).

\section{El Laplaciano Isotrópico}
En \texttt{nucleo\_visual.py}, usamos un stencil (matriz de convolución) de 9 puntos para aproximar el Laplaciano con simetría isotrópica, vital para que los fluidos no parezcan "cuadrados" en la pantalla.
\begin{lstlisting}[language=Python]
# Stencil de 9 puntos (Laplaciano Isotropico ESCALADO espacialmente)
dx_g = 1.2 # Mayor a 1 = Estructuras mas grandes y suaves
lap_u = (0.8 * (U[r+1, c] + U[r-1, c] + U[r, c+1] + U[r, c-1]) + 
         0.2 * (U[r+1, c+1] + U[r+1, c-1] + U[r-1, c+1] + U[r-1, c-1]) - 
         4.0 * U[r, c]) / (dx_g**2)
\end{lstlisting}
El parámetro \texttt{dx\_g = 1.2} expande la grilla computacional, permitiendo formas visuales mucho más orgánicas. Un stencil clásico de 5 puntos produciría artefactos diagonales.

\section{Condiciones de Frontera}
El universo simulado debe tener bordes. Si las partículas o las ondas alcanzan el borde del arreglo bidimensional, ¿qué sucede?
En Cosmic-Generator usamos \textbf{Condiciones de Frontera de Dirichlet Homogéneas} en el Kuramoto-Sivashinsky:
\begin{lstlisting}[language=Python]
out_u[0:2, :] = 0; out_u[-2:, :] = 0
out_u[:, 0:2] = 0; out_u[:, -2:] = 0
\end{lstlisting}
Esto obliga a la turbulencia a morir en los bordes de la pantalla. Por el contrario, en los integradores pseudo-espectrales usamos \textbf{Condiciones Periódicas} implícitas a través de la Transformada Rápida de Fourier (FFT), lo que causa que lo que sale por la izquierda de la pantalla reaparezca por la derecha, como un toroide.



\chapter{Discretizaciones Espaciales}
Para simular el universo continuo en una cuadrícula de píxeles, debemos discretizar los operadores diferenciales continuos (como el Laplaciano $\nabla^2$).

\section{El Laplaciano Isotrópico}
En \texttt{nucleo\_visual.py}, usamos un stencil (matriz de convolución) de 9 puntos para aproximar el Laplaciano con simetría isotrópica, vital para que los fluidos no parezcan "cuadrados" en la pantalla.
\begin{lstlisting}[language=Python]
# Stencil de 9 puntos (Laplaciano Isotropico ESCALADO espacialmente)
dx_g = 1.2 # Mayor a 1 = Estructuras mas grandes y suaves
lap_u = (0.8 * (U[r+1, c] + U[r-1, c] + U[r, c+1] + U[r, c-1]) + 
         0.2 * (U[r+1, c+1] + U[r+1, c-1] + U[r-1, c+1] + U[r-1, c-1]) - 
         4.0 * U[r, c]) / (dx_g**2)
\end{lstlisting}
El parámetro \texttt{dx\_g = 1.2} expande la grilla computacional, permitiendo formas visuales mucho más orgánicas. Un stencil clásico de 5 puntos produciría artefactos diagonales.

\section{Condiciones de Frontera}
El universo simulado debe tener bordes. Si las partículas o las ondas alcanzan el borde del arreglo bidimensional, ¿qué sucede?
En Cosmic-Generator usamos \textbf{Condiciones de Frontera de Dirichlet Homogéneas} en el Kuramoto-Sivashinsky:
\begin{lstlisting}[language=Python]
out_u[0:2, :] = 0; out_u[-2:, :] = 0
out_u[:, 0:2] = 0; out_u[:, -2:] = 0
\end{lstlisting}
Esto obliga a la turbulencia a morir en los bordes de la pantalla. Por el contrario, en los integradores pseudo-espectrales usamos \textbf{Condiciones Periódicas} implícitas a través de la Transformada Rápida de Fourier (FFT), lo que causa que lo que sale por la izquierda de la pantalla reaparezca por la derecha, como un toroide.



\chapter{Seguridad Numérica: Parches Anti-Clipping}
La inyección constante de energía del audio (RMS, graves) puede sobreexcitar el sistema y causar \textbf{Overflows Numéricos (Clipping)} donde los valores explotan hacia el infinito (Not a Number - NaN). Para prevenir esto, los programadores implementaron diversas técnicas de estabilización.

\section{El Parche Tanh en Kuramoto-Sivashinsky}
La ecuación Kuramoto-Sivashinsky es extremadamente inestable debido al término de gradiente al cuadrado no lineal ($-0.5 |\nabla u|^2$). Para evitar su explosión bajo bombos fuertes de audio, aplicamos un "suavizado" asintótico usando la tangente hiperbólica:
\begin{lstlisting}[language=Python]
# Estabilidad KS: Clamping suave con tanh (evita bloqueos numericos)
val = np.tanh(val / 20.0) * 20.0
\end{lstlisting}
Esto restringe físicamente el campo escalar a un rango seguro de $[-20, 20]$, preservando la derivada en el rango lineal (cerca de cero) y frenando el escape exponencial.

\section{Clamping Duro en Reacción-Difusión}
Para los modelos químicos como Gray-Scott, las concentraciones de químicos ($U$ y $V$) no pueden existir físicamente por debajo de 0 o por encima de 1.
\begin{lstlisting}[language=Python]
# Clamping seguro
if val_u > 1.0: val_u = 1.0
elif val_u < 0.0: val_u = 0.0
if val_v > 1.0: val_v = 1.0
elif val_v < 0.0: val_v = 0.0
\end{lstlisting}



\chapter{Seguridad Numérica: Parches Anti-Clipping}
La inyección constante de energía del audio (RMS, graves) puede sobreexcitar el sistema y causar \textbf{Overflows Numéricos (Clipping)} donde los valores explotan hacia el infinito (Not a Number - NaN). Para prevenir esto, los programadores implementaron diversas técnicas de estabilización.

\section{El Parche Tanh en Kuramoto-Sivashinsky}
La ecuación Kuramoto-Sivashinsky es extremadamente inestable debido al término de gradiente al cuadrado no lineal ($-0.5 |\nabla u|^2$). Para evitar su explosión bajo bombos fuertes de audio, aplicamos un "suavizado" asintótico usando la tangente hiperbólica:
\begin{lstlisting}[language=Python]
# Estabilidad KS: Clamping suave con tanh (evita bloqueos numericos)
val = np.tanh(val / 20.0) * 20.0
\end{lstlisting}
Esto restringe físicamente el campo escalar a un rango seguro de $[-20, 20]$, preservando la derivada en el rango lineal (cerca de cero) y frenando el escape exponencial.

\section{Clamping Duro en Reacción-Difusión}
Para los modelos químicos como Gray-Scott, las concentraciones de químicos ($U$ y $V$) no pueden existir físicamente por debajo de 0 o por encima de 1.
\begin{lstlisting}[language=Python]
# Clamping seguro
if val_u > 1.0: val_u = 1.0
elif val_u < 0.0: val_u = 0.0
if val_v > 1.0: val_v = 1.0
elif val_v < 0.0: val_v = 0.0
\end{lstlisting}



\chapter{Análisis del Código: Modo Neural (CPPN)}
A continuación, diseccionamos la red neuronal implementada en NumPy sin frameworks de Deep Learning, maximizando el rendimiento por ciclo de reloj.
\begin{lstlisting}[language=Python]
def _step_lorenz(self, dt=0.01):
    # Resolver inestabilidad de Euler mediante sub-pasos dinamicos
    num_steps = int(max(1, dt / 0.005))
    sub_dt = dt / num_steps
    sigma, rho, beta = 10.0, 28.0, 8.0/3.0
    for _ in range(num_steps):
        x, y, z = self.lorenz_state
        dx = sigma * (y - x) * sub_dt
        dy = (x * (rho - z) - y) * sub_dt
        dz = (x * y - beta * z) * sub_dt
        self.lorenz_state += np.array([dx, dy, dz])
        
        # Clip de seguridad hiper-estricto para evitar Overflow/NaN
        self.lorenz_state = np.clip(self.lorenz_state, -100.0, 100.0)
    return self.lorenz_state / 20.0 # Normalizado a rango latente estable
\end{lstlisting}
Aquí observamos cómo la métrica del audio (\texttt{dt} se escala con \texttt{audio\_z}) acelera o frena el tiempo del universo de Lorenz, y usamos sub-pasos (integración adaptativa manual) para evitar la explosión de Euler si el bombo de la música es demasiado violento.



\chapter{Análisis del Código: Modo Neural (CPPN)}
A continuación, diseccionamos la red neuronal implementada en NumPy sin frameworks de Deep Learning, maximizando el rendimiento por ciclo de reloj.
\begin{lstlisting}[language=Python]
def _step_lorenz(self, dt=0.01):
    # Resolver inestabilidad de Euler mediante sub-pasos dinamicos
    num_steps = int(max(1, dt / 0.005))
    sub_dt = dt / num_steps
    sigma, rho, beta = 10.0, 28.0, 8.0/3.0
    for _ in range(num_steps):
        x, y, z = self.lorenz_state
        dx = sigma * (y - x) * sub_dt
        dy = (x * (rho - z) - y) * sub_dt
        dz = (x * y - beta * z) * sub_dt
        self.lorenz_state += np.array([dx, dy, dz])
        
        # Clip de seguridad hiper-estricto para evitar Overflow/NaN
        self.lorenz_state = np.clip(self.lorenz_state, -100.0, 100.0)
    return self.lorenz_state / 20.0 # Normalizado a rango latente estable
\end{lstlisting}
Aquí observamos cómo la métrica del audio (\texttt{dt} se escala con \texttt{audio\_z}) acelera o frena el tiempo del universo de Lorenz, y usamos sub-pasos (integración adaptativa manual) para evitar la explosión de Euler si el bombo de la música es demasiado violento.



\chapter{Análisis del Código: Modo Neural (CPPN)}
A continuación, diseccionamos la red neuronal implementada en NumPy sin frameworks de Deep Learning, maximizando el rendimiento por ciclo de reloj.
\begin{lstlisting}[language=Python]
def _step_lorenz(self, dt=0.01):
    # Resolver inestabilidad de Euler mediante sub-pasos dinamicos
    num_steps = int(max(1, dt / 0.005))
    sub_dt = dt / num_steps
    sigma, rho, beta = 10.0, 28.0, 8.0/3.0
    for _ in range(num_steps):
        x, y, z = self.lorenz_state
        dx = sigma * (y - x) * sub_dt
        dy = (x * (rho - z) - y) * sub_dt
        dz = (x * y - beta * z) * sub_dt
        self.lorenz_state += np.array([dx, dy, dz])
        
        # Clip de seguridad hiper-estricto para evitar Overflow/NaN
        self.lorenz_state = np.clip(self.lorenz_state, -100.0, 100.0)
    return self.lorenz_state / 20.0 # Normalizado a rango latente estable
\end{lstlisting}
Aquí observamos cómo la métrica del audio (\texttt{dt} se escala con \texttt{audio\_z}) acelera o frena el tiempo del universo de Lorenz, y usamos sub-pasos (integración adaptativa manual) para evitar la explosión de Euler si el bombo de la música es demasiado violento.


\part{Legado: Resumen de la Primera Edición}


\newpage
\part{Arquitectura del Sistema}

\section{Arquitectura del Código y Flujo de Trabajo}

La arquitectura del simulador físico Cosmic-Generator está estructurada para maximizar la modularidad, el rendimiento y la facilidad de uso. El flujo de trabajo o pipeline sigue un modelo secuencial desde la captura de parámetros del usuario hasta la visualización en tiempo real y la síntesis de audio.

\subsection{Punto de Entrada: UI y Configuración (\texttt{app.py})}
El ciclo de ejecución de la plataforma comienza en el archivo \texttt{app.py}. Este script sirve como punto de entrada de la aplicación y utiliza el framework Streamlit para proveer una Interfaz de Usuario (UI) web interactiva. En esta etapa, \texttt{app.py} permite al usuario definir las variables globales y parámetros específicos del sistema (tales como constantes físicas, niveles de fricción, propiedades visuales, entre otros). Al iniciar la simulación, los parámetros se empaquetan y se transmiten al siguiente nivel de la arquitectura.

\subsection{Orquestación del Sistema (\texttt{lanzador.py})}
El módulo \texttt{lanzador.py} cumple el rol de orquestador o controlador principal. Recibe los parámetros configurados en la UI e inicializa todas las piezas fundamentales del simulador. Sus responsabilidades incluyen:
\begin{itemize}
    \item Instanciar el entorno físico.
    \item Configurar los subsistemas de generación de audio y renderizado visual.
    \item Disparar el bucle principal de simulación, delegando la carga computacional a los módulos especializados.
\end{itemize}
De esta forma, se mantiene completamente desacoplada la interfaz web de la lógica pesada del simulador.

\subsection{Patrón Productor-Consumidor en Renderizado (\texttt{src/render/stable/render\_standard.py})}
El rendimiento es crítico en simulaciones físicas de tiempo real. Para asegurar una alta tasa de fotogramas por segundo (FPS), el archivo \texttt{src/render/stable/render\_standard.py} implementa un modelo de concurrencia basado en el patrón \textbf{Productor-Consumidor}, empleando múltiples hilos (threads):

\begin{itemize}
    \item \textbf{Productor (Hilo Físico):} Un hilo en segundo plano se dedica exclusivamente a resolver las ecuaciones matemáticas del simulador. Calcula la cinemática, actualiza posiciones, velocidades y detecta colisiones, produciendo \textit{estados del sistema} que se almacenan de forma segura en una estructura de datos concurrente (cola o buffer).
    \item \textbf{Consumidor (Hilo Gráfico):} El hilo principal de renderizado lee o "consume" los estados físicos ya calculados que están disponibles en el buffer, para dibujar cada fotograma. Al separar estos dos procesos, el cálculo matemático no bloquea el repintado de la pantalla y la visualización gráfica no ralentiza la física, logrando una aceleración notable de los FPS.
\end{itemize}

\subsection{Organización Modular del Código}
El código fuente se organiza de forma rigurosa y cohesiva en distintos directorios bajo \texttt{src/}, respetando el principio de responsabilidad única:

\begin{itemize}
    \item \textbf{\texttt{src/core/}:} Contiene el motor numérico y físico del proyecto. Aquí residen los cálculos de dinámica gravitacional, cinemática y detección de colisiones. Este módulo opera con matemática pura y no tiene conocimiento del entorno gráfico ni sonoro.
    \item \textbf{\texttt{src/audio/}:} Módulo encargado de la síntesis y gestión de sonido. Actúa escuchando los eventos disparados por el Core (por ejemplo, el rebote o impacto de una entidad) y los transforma dinámicamente en efectos de audio espaciales o arreglos musicales.
    \item \textbf{\texttt{src/render/}:} Dedicado de forma exclusiva a la representación visual. Recopila la información geométrica del Core y realiza las llamadas a la biblioteca gráfica pertinente para renderizar el entorno y los elementos simulados en la ventana de la aplicación.
\end{itemize}

\section{Módulo: src/core (Motor Numérico y Físico)}

El directorio \texttt{src/core} constituye el cerebro matemático de la simulación. Su responsabilidad es pura y exclusivamente el cálculo numérico, la dinámica geométrica y la actualización de los estados físicos del universo, permaneciendo completamente ciego respecto a cómo estos datos se dibujan en pantalla o cómo suenan.

\subsection{Archivos Principales y Funcionalidad}
\begin{itemize}
    \item \textbf{\texttt{nucleo\_espectral.py} y \textbf{\texttt{nucleo\_neural.py}}:} Estos archivos contienen las clases matemáticas que resuelven las ecuaciones diferenciales (como Gray-Scott o Navier-Stokes) y las redes neuronales (CPPN). Manejan tensores de NumPy/PyTorch para calcular gradientes espaciales y operadores laplacianos en cada iteración temporal.
    \item \textbf{\texttt{visual\_entities.py}:} Define la orientación espacial, vectores de posición, velocidad y aceleración de las partículas discretas en simulaciones de N-cuerpos. 
    \item \textbf{\texttt{efectos\_visuales.py}:} Gestiona las transformaciones proyectivas de la matriz espacial, tales como las rotaciones en ejes 3D proyectadas a 2D y el factor de escala (zoom) dictaminado por la intensidad física calculada.
\end{itemize}

\subsection{Cálculos de Dinámica y Cinemática}
En este módulo, el tiempo se trata como una variable discreta. En cada actualización (\textit{update}), las coordenadas y matrices de estado evolucionan. Se resuelven integraciones explícitas de Euler o Runge-Kutta para asegurar que la masa, el momento y la energía matemática se comporten bajo leyes deterministas. No existen colores ni píxeles aquí, sólo arreglos de números flotantes de alta precisión.

\section{Módulo: src/audio (Síntesis y Gestión de Sonido)}

El directorio \texttt{src/audio} es el responsable de actuar como el puente sensorial de la simulación. En Cosmic-Generator, el flujo de información es bidireccional: la música altera la física, pero la física (eventos) también debe ser interpretada.

\subsection{Extracción Acústica (\texttt{audio\_analyzer.py})}
Este archivo emplea la librería \texttt{librosa} para diseccionar un archivo de audio de entrada en variables numéricas utilizables por el motor físico. Extrae:
\begin{itemize}
    \item \textbf{Energía (RMS):} La amplitud o volumen del fotograma actual.
    \item \textbf{Spectral Flux (Flujo Espectral):} Los ataques transitorios (ej: golpes de batería o platillos).
    \item \textbf{Tonalidad y Pitch Centroide:} Las frecuencias fundamentales que definirán las paletas cromáticas (\textit{Cromestesia}).
\end{itemize}

\subsection{Alineación de Letras (\texttt{motor\_lyrics.py})}
Contiene el motor de Inteligencia Artificial (\texttt{stable-ts} / Whisper). Escucha el archivo de audio, aísla las frecuencias vocales y genera un diccionario JSON con las marcas de tiempo (timestamps) exactas de cada palabra. Además, implementa la tecnología de \textbf{Alineación Forzada}, permitiendo al usuario inyectar la letra original y obligando a la IA a sincronizar matemáticamente las sílabas, eliminando cualquier tipo de alucinación del modelo.

\subsection{Eventos Disparados por el Core}
Cuando el módulo físico (\texttt{src/core}) detecta una colisión de alta energía, envía una señal asíncrona a este módulo. Dependiendo de la magnitud del impacto, el sistema de audio puede inyectar un efecto sonoro o alterar la retroalimentación visual basada en frecuencias sonoras inyectadas de vuelta al pipeline.

\section{Módulo: src/render (Representación Visual Exclusiva)}

El directorio \texttt{src/render} es el final de la cadena de montaje computacional (Pipeline). Es el módulo consumidor puro: no realiza cálculos físicos ni análisis de audio; su única misión es recopilar los tensores numéricos del \texttt{core}, asignarles texturas y enviarlos a la pantalla o a un archivo de video final de manera ultarrápida.

\subsection{Estructura del Motor de Renderizado}
El directorio se divide en ramas de estabilidad:
\begin{itemize}
    \item \textbf{\texttt{stable/render\_standard.py}:} El bucle principal de renderizado usando OpenCV (\texttt{cv2}). Se encarga de transformar las matrices abstractas del \texttt{core} en fotogramas RGB (\textit{frames}).
    \item \textbf{\texttt{experimental/}:} Contiene motores de renderizado más avanzados o no lineales (como trazado de rayos básico o Shaders nativos) que están en fase beta.
\end{itemize}

\subsection{Proceso de Dibujado (Drawing Pipeline)}
El motor \texttt{render\_standard.py} orquesta la función \texttt{generar\_animacion\_god\_mode()}. Su ciclo de vida es el siguiente:
\begin{enumerate}
    \item Lee la matriz escalar generada por \texttt{src/core} (por ejemplo, el mapa de densidad de un fluido).
    \item Aplica un mapeo de color (Colormap) dinámico basado en las variables dictadas por \texttt{src/audio} (Tonalidad).
    \item Si corresponde, dibuja las partículas usando operaciones de trazado matricial o \texttt{cv2.circle}.
    \item Si el \texttt{motor\_lyrics.py} envía texto sincronizado, invoca rutinas tipográficas (\texttt{Pillow} o \texttt{cv2.putText}) para sobreimprimir la lírica en pantalla con colores vibrantes y desenfoques (Bloom).
    \item Escribe el fotograma directamente en el \textit{VideoWriter} (FFmpeg) para exportar el `.mp4` final, reportando el progreso a la interfaz de usuario.
\end{enumerate}

\part{El Puente Sensorial (Audio a Física)}

\section{Extracción de Parámetros de Audio}
\label{sec:extraccion_audio}

El corazón rítmico de Cosmic-Generator reside en su módulo de análisis avanzado, ubicado en \texttt{src/audio/motor\_audio.py}. Para lograr una reactividad visual ultra optimizada y matemáticamente precisa, empleamos la librería \texttt{librosa} para diseccionar el espectro acústico en el dominio del tiempo y la frecuencia, transformándolo en tensores de control que alimentan directamente nuestra tubería gráfica.

\subsection{Energía y Zoom (RMS)}
El cálculo del \textit{Root Mean Square} (RMS) nos proporciona una métrica directa de la energía térmica y amplitud de la señal.
Extraemos la energía promedio por cada \textit{frame} de audio y la aplicamos como un escalar no lineal en la matriz de transformación proyectiva (Zoom). De esta manera, a mayor intensidad percusiva o volumen acústico, generamos una compresión espacial que empuja la cámara hacia la acción.

\subsection{Turbulencia y Flujo Espectral (Spectral Flux)}
Para inducir caos controlado e inestabilidades realistas en la simulación de partículas y fluidos, calculamos el \textit{Spectral Flux}. Esta métrica evalúa la tasa de cambio (la primera derivada discreta) de la magnitud del espectro de potencias entre \textit{frames} consecutivos.
Este gradiente espectral alimenta de inmediato nuestros campos vectoriales y los coeficientes de advección, inyectando vórtices y ráfagas de turbulencia direccional proporcionales a los ataques e impulsos de alta frecuencia de la pista.

\subsection{Cromestesia y Color (Pitch/Tonalidad)}
Nuestra arquitectura de sombreadores de fragmentos (\textit{Fragment Shaders}) no emplea paletas de colores estáticas. El color de la simulación obedece a un mapeo dinámico entre el \textit{pitch} dominante y el espacio cromático.
Mediante \texttt{librosa}, extraemos la frecuencia fundamental ($f_0$) y el cromagrama del audio para clasificar la tonalidad en tiempo real. Esta tonalidad modula los canales RGB de forma algorítmica, induciendo una cromestesia donde las frecuencias graves tiñen el entorno de fluidos densos y rojizos, mientras que los armónicos agudos cristalizan la escena en destellos cianes y blancos de alta luminancia.

\section{Frecuencia y Sincronización de Cambio de Parámetros}
\label{sec:frecuencia_parametros}

La sincronización milimétrica entre la integración del motor físico en la GPU y la línea de tiempo acústica es nuestro mayor desafío resuelto. Hemos estructurado una arquitectura asíncrona altamente paralela para evitar bloqueos y asegurar que el renderizado se mantenga fluido sin sacrificar la rigurosidad de las colisiones y campos de fuerza.

\subsection{Muestreo Discreto del Audio}
En la fase de pre-cómputo, el módulo de audio procesa la pista completa extrayendo características \textit{frame} a \textit{frame}. Estableciendo una frecuencia de muestreo de fotogramas (por ejemplo, a 30 FPS constantes), dividimos el audio obteniendo un arreglo vectorial de tamaño $\text{duración total} \times \text{FPS}$.
Este arreglo indexable nos permite un acceso a memoria en $\mathcal{O}(1)$ para extraer las variables ambientales necesarias de acuerdo al tiempo de la simulación.

\subsection{Time-stepping Loop vs. Render Loop}
Para maximizar el rendimiento gráfico, implementamos un desacoplamiento de bucles.
El bucle de dibujado (Render Loop) está dedicado exclusivamente a interpolar y enviar \textit{buffers} a la pantalla con la menor latencia posible.
Simultáneamente, en un hilo de ejecución paralelo, transcurre el \textit{Time-stepping Loop} de la simulación física. En CADA iteración de este bucle físico, el motor extrae el \textit{timestamp} actual e indexa el \textit{array} de audio, inyectando los parámetros del \textit{frame} acústico exactamente en el momento que la física requiere su actualización.

\subsection{Mutación Dinámica de la Difusión}
El verdadero poder del Cosmic-Generator se despliega en cómo la acústica altera las leyes físicas. En cada actualización de los \textit{solvers} de la simulación de fluidos, los coeficientes de las ecuaciones de difusión mutan drásticamente:
la intensidad acústica (derivada del RMS) modula escalarmente la viscosidad y difusividad térmica, forzando a los fluidos a volverse más reactivos o densos.
Simultáneamente, las frecuencias dominantes inyectan anisotropía direccional al tensor de difusión. El resultado es un estado de materia puramente reactivo, donde el sonido altera topológicamente los coeficientes matemáticos a razón de 30 veces por segundo, esculpiendo gráficamente la canción.

\section{El Puente Físico-Sonoro: Destruyendo la Caja Negra}

Este documento revela el mecanismo exacto, a nivel matemático y de código, mediante el cual las señales de audio se inyectan en los cinco motores físicos del Cosmic-Generator. No hay magia ni cajas negras; hay un mapeo directo y determinista entre el análisis de la señal y las ecuaciones diferenciales parciales.

\subsection{Extracción de Parámetros de Audio}
Para cada fotograma (típicamente a 30 o 60 FPS), el sistema procesa una ventana de la señal de audio de entrada y extrae un vector de características espectrales y temporales. Estas variables de control se normalizan y se aplican filtros de suavizado antes de integrarse en los métodos numéricos.

\begin{itemize}
    \item \textbf{Amplitud (RMS) -- $A(t)$:} Representa la energía global instantánea. Controla variables macroscópicas, inyecciones de energía de fondo y tasas de alimentación.
    \item \textbf{Flujo Espectral (Spectral Flux) -- $S(t)$:} Mide los cambios abruptos de energía entre frames (transitorios, golpes de percusión). Se usa para inyectar perturbaciones impulsivas y forzamientos abruptos.
    \item \textbf{Tonalidad / Centroide Espectral -- $T(t)$:} Representa el `color' o brillo del sonido. Modula propiedades intrínsecas del medio, como coeficientes de difusión, repulsión o velocidad de onda.
\end{itemize}

\subsection{Motor 1: Reacción-Difusión (Gray-Scott / Corrosión)}
El motor Gray-Scott simula la interacción química de dos sustancias, $u$ y $v$, generando laberintos, células y patrones de Turing. Las ecuaciones que evolucionamos son:
\begin{align*}
\frac{\partial u}{\partial t} &= D_u \nabla^2 u - u v^2 + f(t)(1 - u) \\
\frac{\partial v}{\partial t} &= D_v \nabla^2 v + u v^2 - (f(t) + k(t)) v
\end{align*}

\textbf{Inyección de Audio:}
Las constantes $f$ y $k$, que en la formulación original dictan la `especie' del patrón, mutan dinámicamente con el audio:
\begin{itemize}
    \item \textbf{Tasa de Alimentación (Feed Rate, $f$):} La amplitud $A(t)$ incrementa la tasa $f$. Un sonido denso introduce más químico $u$, hinchando los patrones hasta volverlos expansivos e interconectados.
    $$f(t) = f_{base} + \alpha_f A(t)$$
    \item \textbf{Tasa de Mortalidad (Kill Rate, $k$):} Se perturba negativamente con los golpes rítmicos ($S(t)$). Un `kick' disminuye momentáneamente $k$, provocando explosiones celulares rápidas.
    $$k(t) = k_{base} - \alpha_k S(t)$$
\end{itemize}

\subsection{Motor 2: Kuramoto-Sivashinsky (Fuego / Caos Espacial)}
Este motor modela inestabilidades de propagación de frentes (como llamas) exhibiendo turbulencia espacio-temporal. La ecuación escalar modificada es:
\begin{equation}
\frac{\partial u}{\partial t} + \nu \nabla^4 u + \mu(t) \nabla^2 u + \frac{1}{2}|\nabla u|^2 = F(x, y, t)
\end{equation}

\textbf{Inyección de Audio:}
\begin{itemize}
    \item \textbf{Coeficiente de Inestabilidad ($\mu$):} La amplitud o intensidad de los graves aumenta el coeficiente $\mu$, inyectando energía en longitudes de onda largas y provocando que el `fuego' se vuelva más turbulento y desordenado.
    $$\mu(t) = \mu_0 + \alpha_\mu A(t)$$
    \item \textbf{Forzamiento Espacial ($F$):} Los transitorios $S(t)$ activan mapas de ruido espacial en instantes puntuales, actuando como inyecciones o `chispas' que alimentan el campo escalar localizado, perturbando la ecuación bruscamente.
\end{itemize}

\subsection{Motor 3: Gross-Pitaevskii (Cuántico / Superfluido)}
Describe el comportamiento microscópico de un condensado de Bose-Einstein o de un superfluido, modelado sobre una función de onda compleja $\psi$.
\begin{equation}
i\hbar \frac{\partial \psi}{\partial t} = \left( -\frac{\hbar^2}{2m}\nabla^2 + V(\mathbf{x}, t) + g(t)|\psi|^2 \right)\psi
\end{equation}

\textbf{Inyección de Audio:}
\begin{itemize}
    \item \textbf{Potencial Externo ($V$):} El audio genera un campo topológico sobre el lienzo. La intensidad global $A(t)$ modula la profundidad o barreras del potencial atrapando o liberando el fluido cuántico en zonas de alta densidad energética visual.
    \item \textbf{Constante de Auto-Interacción ($g$):} Se asocia a la Tonalidad o Centroide $T(t)$. Si el sonido tiene frecuencias predominantemente agudas, el coeficiente $g$ se vuelve más repulsivo (o menos atractivo), alterando el tamaño y la velocidad de nucleación de los vórtices cuánticos.
\end{itemize}

\subsection{Motor 4: Onda Clásica}
Simula la propagación de deformaciones transversales en una malla 2D ideal (análogo al parche de un tambor).
\begin{equation}
\frac{\partial^2 u}{\partial t^2} = c(t)^2 \nabla^2 u - \gamma(t) \frac{\partial u}{\partial t} + P(\mathbf{x}, t)
\end{equation}

\textbf{Inyección de Audio:}
\begin{itemize}
    \item \textbf{Velocidad de Propagación ($c$):} Se vuelve una función dinámica ligada a la frecuencia fundamental del audio. Sonidos más agudos `tensan' el medio virtual, acelerando la velocidad de fase $c$.
    \item \textbf{Amortiguamiento ($\gamma$):} La amplitud sostenida reduce el rozamiento interno. Un drone musical provoca que las ondas perduren en el espacio; el silencio incrementa el amortiguamiento limpiando el lienzo.
    \item \textbf{Excitación ($P$):} Golpes rítmicos mapean fuerza pura en coordenadas específicas, emitiendo verdaderas frentes de onda en sincronía.
\end{itemize}

\subsection{Motor 5: Korteweg-de Vries (KdV / Tsunamis, Solitones)}
La ecuación KdV (o su extensión 2D Kadomtsev-Petviashvili) permite la emergencia de solitones, donde el empinamiento no lineal balancea exactamente la dispersión de las ondas.
\begin{equation}
\frac{\partial u}{\partial t} + \alpha(t) u \frac{\partial u}{\partial x} + \beta(t) \frac{\partial^3 u}{\partial x^3} = 0
\end{equation}

\textbf{Inyección de Audio:}
\begin{itemize}
    \item \textbf{Efecto No Lineal ($\alpha$):} Está directamente controlado por la energía RMS de la señal de audio. Al incrementarse el volumen general, $\alpha$ crece desproporcionadamente, provocando que los frentes de onda se empinen formando picos afilados o estructuras tipo `tsunami'.
    $$\alpha(t) = \alpha_{base} \cdot (1 + \lambda_A A(t))$$
    \item \textbf{Efecto Dispersivo ($\beta$):} Responde al perfil espectral $T(t)$. Frecuencias muy agudas y brillantes aumentan el término dispersivo diferencial, dividiendo grandes frentes de onda en un tren múltiple de solitones de alta frecuencia.
\end{itemize}

\vspace{0.5cm}
\textbf{Conclusión:}
En Cosmic-Generator, el audio no actúa como un mero filtro visual en etapa de post-procesamiento. Las propiedades rítmicas, dinámicas y armónicas de la música se convierten en perturbaciones, coeficientes no lineales y tasas de propagación insertadas directamente en los integradores Euler, Runge-Kutta de 4to orden y esquemas de Diferencias Finitas (FDTD). La música rediseña las leyes físicas del universo por cada $\Delta t$.

\part{Motores Físicos y Ecuaciones Diferenciales}

\section{Fundamentos Físico-Matemáticos de los Motores Dinámicos}
\label{sec:fundamentos_motores}

En esta sección, formalizamos los modelos fenomenológicos que subyacen a los diversos motores visuales del simulador \textit{Cosmic-Generator}. La rigurosidad nos exige comprender que los patrones observados no son meros artefactos gráficos, sino soluciones emergentes de sistemas de Ecuaciones Diferenciales Parciales (EDP) altamente no lineales.

\subsection{Motor de Corrosión: Modelo de Reacción-Difusión de Gray-Scott}
El motor que produce los intrincados patrones biológicos y de "corrosión" se basa en el célebre sistema de reacción-difusión introducido por P. Gray y S.K. Scott. Este modelo describe la cinética de dos especies químicas, $U$ y $V$, que se difunden en el espacio y reaccionan según una dinámica de retroalimentación autocatalítica. El sistema acoplado se rige por:

\begin{equation}
\frac{\partial u}{\partial t} = D_u \nabla^2 u - u v^2 + f(1 - u)
\end{equation}

\begin{equation}
\frac{\partial v}{\partial t} = D_v \nabla^2 v + u v^2 - (f + k)v
\end{equation}

donde $u(\mathbf{r},t)$ y $v(\mathbf{r},t)$ representan las concentraciones espaciales. $D_u$ y $D_v$ son las constantes de difusión, $f$ es la tasa de alimentación (feed) que repone $U$, y $k$ la tasa de decaimiento o aniquilación (kill) de $V$. Los patrones visuales emergen de la inestabilidad de Turing inducida por la competencia entre una difusión rápida de la especie inhibidora y la reacción local de la especie activadora.

\subsection{Motor de Fuego: Ecuación de Kuramoto-Sivashinsky}
Para la simulación de frentes de llama y turbulencia espaciotemporal, empleamos la ecuación de Kuramoto-Sivashinsky. Originalmente derivada para modelar inestabilidades termo-difusivas en frentes de combustión, su forma fenomenológica en términos de fluctuaciones del frente $h(\mathbf{r},t)$ es:

\begin{equation}
\frac{\partial h}{\partial t} = -\frac{1}{2} |\nabla h|^2 - \nabla^2 h - \nabla^4 h
\end{equation}

Aquí observamos tres contribuciones fundamentales: el término no lineal $\frac{1}{2} |\nabla h|^2$ que estabiliza el crecimiento y transfiere energía entre escalas; el término difusivo inverso $-\nabla^2 h$ que bombea energía al sistema (inestabilidad en bajas frecuencias); y el término hiperdifusivo $-\nabla^4 h$ que disipa las perturbaciones de onda corta. Esto genera el comportamiento caótico y fractal propio del "fuego".

\subsection{Motor Cuántico: Ecuación de Gross-Pitaevskii}
Para representar campos de naturaleza cuántica, como los condensados de Bose-Einstein o fluidos de luz, recurrimos a una formulación de campo medio que obedece a la Ecuación de Schrödinger No Lineal (NLS), conocida en materia condensada como ecuación de Gross-Pitaevskii:

\begin{equation}
i \hbar \frac{\partial \psi}{\partial t} = \left( -\frac{\hbar^2}{2m} \nabla^2 + V(\mathbf{r}) + g |\psi|^2 \right) \psi
\end{equation}

donde $\psi(\mathbf{r},t)$ es la función de onda macroscópica compleja, $m$ es la masa efectiva, $V(\mathbf{r})$ es un potencial externo de confinamiento, y el término $g |\psi|^2$ representa las interacciones de contacto entre partículas (repulsivas para $g > 0$). Las fluctuaciones en su amplitud y fase dan origen a vórtices cuánticos y a una interferencia dinámica soberbia.

\subsection{Motor de Onda: Dinámica Hiperbólica Clásica}
El motor de ondulaciones superficiales es, en el fondo, una integración de la ecuación de onda clásica, el arquetipo de las ecuaciones diferenciales hiperbólicas:

\begin{equation}
\frac{\partial^2 u}{\partial t^2} = c^2 \nabla^2 u - \gamma \frac{\partial u}{\partial t}
\end{equation}

El campo escalar $u(\mathbf{r},t)$ representa la elongación respecto al equilibrio. El parámetro $c$ denota la velocidad de propagación de la perturbación, y hemos añadido un término fenomenológico $\gamma \partial_t u$ para modelar la disipación (amortiguamiento). La rigurosa conservación de la energía y el principio de Huygens se manifiestan visualmente a través de la reflexión y refracción.

\subsection{Motor de Tsunamis: Ecuación de Korteweg-de Vries (KdV)}
Finalmente, la asombrosa persistencia de ondas solitarias en el entorno visual está comandada por la ecuación de Korteweg-de Vries (KdV). Modelando aguas poco profundas, nos entrega la génesis de los \textit{solitones}:

\begin{equation}
\frac{\partial u}{\partial t} + \alpha u \frac{\partial u}{\partial x} + \beta \frac{\partial^3 u}{\partial x^3} = 0
\end{equation}

Esta maravilla matemática describe cómo la dispersión lineal (el término con tercera derivada $\partial^3_x u$) y el empinamiento no lineal (el término advectivo $u \partial_x u$) se balancean con exactitud. El resultado es un paquete de ondas que se propaga sin deformarse temporalmente, un fenómeno crítico para la descripción de los tsunamis y ciertas fibras ópticas.

\section{Discretización y Simulación Numérica de Sistemas Dinámicos}
\label{sec:simulacion_numerica}

Es un principio innegable en la física computacional moderna que una Ecuación Diferencial Parcial carece de utilidad empírica hasta que no es llevada a un entorno discreto estable. En \textit{Cosmic-Generator}, materializamos las teorías del Capítulo 1 utilizando tensores en PyTorch o arreglos contiguos en NumPy, acelerados frecuentemente mediante el compilador JIT Numba. El éxito de estos métodos recae en la aplicación canónica de Diferencias Finitas y algoritmos de integración temporal sólidos.

\subsection{Discretización Espacial y el Operador Laplaciano}
El núcleo de la mayoría de nuestros motores (Gray-Scott, Onda, Kuramoto-Sivashinsky) depende de la evaluación precisa del Laplaciano espacial $\nabla^2 u$. En un retículo bidimensional de espaciado uniforme $\Delta x = \Delta y$, aplicamos un esquema de diferencias finitas de segundo orden (la plantilla de 5 puntos de Von Neumann):

\begin{equation}
\nabla^2 u_{i,j} \approx \frac{u_{i+1,j} + u_{i-1,j} + u_{i,j+1} + u_{i,j-1} - 4u_{i,j}}{(\Delta x)^2}
\end{equation}

Para las evaluaciones dispersivas o de mayor orden, como el término bilaplaciano $\nabla^4 h$ en Kuramoto-Sivashinsky o la tercera derivada en KdV, se componen iterativamente estas aproximaciones o se aplican plantillas extendidas de 9 o más puntos para preservar la simetría rotacional en la grilla y evitar artefactos anisotrópicos puramente computacionales.

\subsection{Condiciones de Frontera Periódicas y Optimizaciones Tensoriales}
Un simulador continuo precisa una topología sin bordes para evitar reflexiones espurias, asumiendo una geometría toroidal. Esto lo logramos matemáticamente imponiendo condiciones de frontera periódicas: $u(0, y) = u(L_x, y)$ y $u(x, 0) = u(x, L_y)$.

A nivel de software (NumPy o PyTorch), esto se vectoriza magistralmente mediante la operación \texttt{roll}. En lugar de utilizar lentos bucles \texttt{for}, desplazamos el estado global del sistema de esta manera:

\begin{equation}
u_{i+1,j} \implies \texttt{np.roll(u, shift=-1, axis=0)}
\end{equation}

Sumar matrices desplazadas permite computar las diferencias finitas sobre toda la malla simultáneamente, sacando provecho de las instrucciones SIMD del procesador o la paralelización extrema de las arquitecturas CUDA en la GPU.

\subsection{Integración Temporal: Euler Explícito vs. Runge-Kutta}
Toda vez calculado el funcional de derivadas espaciales $F(u)$, la evolución se transforma en un problema de valor inicial: $\partial_t u = F(u)$. Para los motores que demandan máximo rendimiento en tiempo real y exhiben alta disipación (como Gray-Scott o la Onda clásica), utilizamos el método de \textbf{Euler explícito}, el cual aproxima la serie de Taylor al primer orden:

\begin{equation}
u(t + \Delta t) = u(t) + \Delta t \cdot F(u(t))
\end{equation}

Aunque computacionalmente económico, el método de Euler exige un paso temporal $\Delta t$ severamente restringido por el criterio de estabilidad de Courant-Friedrichs-Lewy (CFL), especialmente ante los términos difusivos ($\Delta t \propto (\Delta x)^2$). 

Por el contrario, para ecuaciones rígidas ("stiff") y sin disipación natural como la Ecuación de Schrödinger (Gross-Pitaevskii) o la KdV, el método de Euler introduce una espuria explosión de energía. En estas dinámicas, el simulador debe emplear esquemas simplécticos o el clásico \textbf{Runge-Kutta de cuarto orden (RK4)}, integrando estados intermedios $k_1, k_2, k_3, k_4$ que cancelan errores hasta el orden $\mathcal{O}(\Delta t^4)$, asegurando la conservación de los invariantes físicos fundamentales que la academia manda.

\section{Fundamentación de los Parámetros Físicos Base}

Para que Cosmic-Generator produzca un espectáculo visual estable y estéticamente superior antes de que la acústica intervenga, fue imperativo seleccionar un conjunto de parámetros físicos iniciales (base) ubicados en regímenes matemáticos muy específicos. A continuación, el Comité de Física Teórica desglosa la elección de estos escalares codificados en \texttt{config.py}.

\subsection{Motor Gray-Scott (Reacción-Difusión)}
El motor principal del simulador arranca frecuentemente con los parámetros:
\begin{equation}
f \text{ (Feed Rate)} = 0.055, \quad k \text{ (Kill Rate)} = 0.062
\end{equation}

\textbf{¿Por qué?} Estos valores exactos fueron descubiertos por John E. Pearson en 1993 ("Complex Patterns in a Simple System"). En el diagrama de fase paramétrico $(f, k)$ del modelo de Gray-Scott, esta coordenada corresponde al régimen de \textbf{Patrones de Coral o Laberintos}. 
Si el \textit{feed rate} fuese más alto, el sistema colapsaría a un estado homogéneo. Si fuese más bajo, la reacción se extinguiría. Este punto crítico permite que, cuando el audio introduzca fluctuaciones vía RMS, el laberinto se fragmente transitoriamente en manchas de Turing (spots) o mitosis celular, retornando al laberinto al cesar la música.

\subsection{Motor Kuramoto-Sivashinsky (KS)}
Los parámetros base dictados para el fuego turbulento son:
\begin{equation}
\nu \text{ (Hiperdifusión)} = 0.1, \quad \text{Advección} = 0.0
\end{equation}

\textbf{¿Por qué?} La ecuación KS es un modelo arquetípico de caos espaciotemporal. Al fijar el coeficiente de hiperdifusión ($\nabla^4$) en 0.1, garantizamos que las inestabilidades de onda larga generadas por la difusión inversa ($-\nabla^2$) sean amortiguadas exactamente en el rango de los píxeles de nuestra resolución. Esto produce celdas de "fuego" que tienen un diámetro visualmente armónico en resoluciones 720p/1080p, en lugar de ruido blanco de alta frecuencia que resultaría desagradable a la vista.

\subsection{Motor Cuántico (Gross-Pitaevskii - GPE)}
Para la simulación del condensado de Bose-Einstein, el parámetro de interacción $g$ está seteado en:
\begin{equation}
g \text{ (Interacción)} = -2.0
\end{equation}

\textbf{¿Por qué?} En la física de gases cuánticos diluidos, un valor de $g < 0$ indica interacciones \textbf{atractivas} entre bosones. Matemáticamente, esto induce el colapso del condensado y la formación de solitones brillantes (Bright Solitons). Visualmente, esto crea un efecto de "agujero negro luminoso" donde la materia colapsa estéticamente hacia el centro, el cual luego es perturbado y puesto en rotación por los ataques percusivos de la batería.

\subsection{Motor de Onda Clásica}
El amortiguamiento ($\gamma$) se establece muy cercano a $0.99$ (es decir, decaimiento del 1\%):
\begin{equation}
\gamma = 0.99
\end{equation}

\textbf{¿Por qué?} Un amortiguamiento nulo ($\gamma = 1.0$) causaría que la energía acústica inyectada fotograma a fotograma se acumule infinitamente, sobreexponiendo la pantalla a blanco (clipping). Un decaimiento del 1\% provee un balance perfecto: las ondas persisten el tiempo suficiente para generar interferencias (patrones de Moiré) al chocar con las paredes periódicas, pero se disipan justo a tiempo para dejar espacio visual al próximo transitorio de la canción.

\part{Efectos Especiales y Geometría Sagrada (Motores Inéditos)}


\chapter{Sistema de Partículas: Los Espíritus Procedurales}
Los "Espíritus" son entidades volumétricas reactivas. El sistema de N-cuerpos en \texttt{visual\_entities.py} simula humo browniano.
Las partículas son inyectadas basándose en el parámetro acústico \textbf{Kick} (Bombo) y su tiempo de vida es alargado por el contenido armónico.

\section{Código de Emisión (Extracto)}
\begin{lstlisting}[language=Python]
# EMITIR NUEVAS PARTICULAS (Humo volumetrico)
emision_rate = int(30 + kick * 100) # Mas particulas si hay kick

# Velocidad: Movimiento organico ascendente/aleatorio
base_vx = np.sin(t * 2 + self.noise_offset) * 2.0
base_vy = -2.0 - kick * 3.0 # Hacia arriba (gravedad invertida)
\end{lstlisting}
El humo es renderizado mediante \textit{additive blending} con \texttt{cv2.circle}, generando volúmenes suaves en tiempo real.

\chapter{Gestión Lumínica: Bloom y Flash}
El motor evita procesadores gráficos complejos aislando altas luces matemáticamente mediante pirámides de downsampling iterativo, creando el efecto \textbf{Bloom} o neón.

\section{Código del Multi-Layer Glow (Extracto)}
\begin{lstlisting}[language=Python]
def aplicar_bloom(self, img, intensity, threshold=200):
    # 1. Extraer altas luces
    gray = cv2.cvtColor(img, cv2.COLOR_BGR2GRAY)
    _, mask = cv2.threshold(gray, threshold, 255, cv2.THRESH_BINARY)
    bright = cv2.bitwise_and(img, img, mask=mask)
    
    # 2. Multi-Layer Glow (Piramide Gaussiana)
    h, w = bright.shape[:2]
    w1, h1 = max(w // 2, 1), max(h // 2, 1)
    l1 = cv2.resize(bright, (w1, h1), interpolation=cv2.INTER_LINEAR)
    b1 = cv2.GaussianBlur(l1, (5, 5), 0)
    # [...]
\end{lstlisting}

\chapter{Extracción Lirical y Alineación Forzada}
Usando Whisper (\texttt{stable-ts}), el código aisla los fonemas en el archivo \texttt{motor\_lyrics.py}. Para evitar la "Caja Negra" de la IA, usamos \textit{Forced Alignment}, que mapea la letra oficial a marcas de tiempo.

\begin{lstlisting}[language=Python]
if external_text:
    print("🎯 MODO ALINEACION: Usando letra original para sincronizar...")
    # align() fuerza al modelo a usar este texto y solo buscar los tiempos
    result = model.align(
        audio_array, 
        external_text, 
        language='es',
        progress_callback=self.progress_callback
    )
\end{lstlisting}
Esto se envía al renderizador para sobreimprimir texto estilizado al compás exacto de la voz del vocalista.

\chapter{Geometría Sagrada y Caos 3D}
Las ilusiones fractales y la simetría sagrada (efecto Caleidoscopio) se implementan invirtiendo y superponiendo tensores, mientras que los atractores de Lorenz (Caos 3D) se proyectan en el plano bidimensional modificando su matriz de rotación con los transitorios del sonido.

\begin{lstlisting}[language=Python]
def kaleidoscope(frame, splits=4):
    h, w = frame.shape[:2]
    # Cortar un fragmento y reflejarlo simetricamente
    # para producir una geometria de loto / fractal
\end{lstlisting}
Esto inyecta espiritualidad geométrica directamente correlacionada con las frecuencias perfectas de las octavas musicales (Tonalidad Pura).


\end{document}
